\section{半胞 chiral 对称}
\label{sec:chiral}

有限 Bloch 直和降低了整体维数，但通常不会继续简化单个 fiber。当半周期平移反转
step-two 系数 word 时，反转 step-one 部分所需的交错符号会产生一个反对易单项式算子。
下面的系数计算在这一自然算子类中完整刻画该机制。

\subsection{半胞算子}

令 $p=2m$，并在 lift 上定义
\begin{equation}
 (Dx)_i=(-1)^ix_i,\qquad (T_mx)_i=x_{i+m},\qquad K_m=DT_m.
 \label{eq:half-cell-operators}
\end{equation}
$D$ 反转所有位移一耦合并保持位移二耦合，$T_m$ 比较两个半胞，因此 $K_m$ 是最自然的
chiral 候选。

\begin{theorem}
\label{thm:coefficient-chiral}
对 $2m$-周期符号 word $\tau$，以下两项等价：
\begin{enumerate}
 \item $K_mA_\tau=-A_\tau K_m$；
 \item 对每个 $i\in\mathbb Z$，$\tau_{i+m}=-\tau_i$。
\end{enumerate}
\end{theorem}

\begin{proof}
把反对易子作用到任意序列，位移一项相消，剩余系数为
\begin{align}
 ((A_\tau K_m+K_mA_\tau)x)_i=(-1)^i\bigl[&
 (\tau_{i-2}+\tau_{i+m-2})x_{i+m-2}\notag\\
 &+(\tau_i+\tau_{i+m})x_{i+m+2}\bigr].
 \label{eq:anticommutator-coefficients}
\end{align}
后半 word 是前半 word 的相反数时两项均消失。反之，若反对易子对所有序列为零，则
$x_{i+m+2}$ 的系数为零，从而 $\tau_{i+m}=-\tau_i$。
\end{proof}

$K_m$ 保持每个 $\mathcal B_z^{(2m)}$，但尚未归一化为对合。在该空间上
\begin{equation}
 T_m^2=zI,\qquad T_mD=(-1)^mDT_m,\qquad K_m^2=(-1)^mzI.
 \label{eq:km-square}
\end{equation}
选择单位复数 $\gamma_m(z)$ 使
\begin{equation}
 \gamma_m(z)^2=(-1)^mz^{-1}.
 \label{eq:gamma-normalization}
\end{equation}
若 $\xi^2=z$，则 $m$ 为偶数时可取 $\xi^{-1}$，$m$ 为奇数时可取
$\mathrm i\xi^{-1}$。于是
\begin{equation}
 J_z=\gamma_m(z)K_m|_{\mathcal B_z^{(2m)}}
 \label{eq:normalized-j}
\end{equation}
是酉对合，因而自伴。改变 $z$ 的平方根只把 $J_z$ 换成 $-J_z$，交换两个 chiral 子空间。

\subsection{Flux 判据}

\begin{theorem}
\label{thm:flux-chiral}
对 $2m$-周期符号 word，以下两项等价：
\begin{enumerate}
 \item 对所有 $i$，$\tau_{i+m}=-\tau_i$；
 \item 对所有 $i$，$Q_{i+m}=Q_i$，且 $\prod_{j=0}^{m-1}Q_j=-1$。
\end{enumerate}
此时
\begin{equation}
 J_zH_\tau^{(2m)}(z)=-H_\tau^{(2m)}(z)J_z.
 \label{eq:fiber-anticommutation}
\end{equation}
\end{theorem}

\begin{proof}
若 $\tau_{i+m}=-\tau_i$，则
$Q_{i+m}=(-\tau_i)(-\tau_{i+1})=Q_i$，并由望远镜相消得到
$\prod_{j=0}^{m-1}Q_j=\tau_0\tau_m=-1$。

反之，令 $r_i=\tau_{i+m}/\tau_i$。$Q$ 的半周期性给出 $r_ir_{i+1}=1$，故所有
$r_i$ 等于同一常数 $r$。再次望远镜相消可得
$r=\tau_m/\tau_0=\prod_{j=0}^{m-1}Q_j=-1$。最后由
定理~\ref{thm:coefficient-chiral} 得到反对易关系。
\end{proof}

\subsection{代数后果}

$J_z$ 的单项式矩阵没有固定坐标，故迹为零。它在 $2m$ 维空间上的两个本征子空间维数均为
$m$，反对易关系把二者相互映射。在适应该分解的酉基下，
\begin{equation}
 H_\tau^{(2m)}(z)\simeq
 \begin{pmatrix}0&B(z)^*\\B(z)&0\end{pmatrix},\qquad
 H_\tau^{(2m)}(z)^2\simeq
 \begin{pmatrix}B(z)^*B(z)&0\\0&B(z)B(z)^*\end{pmatrix}.
 \label{eq:general-chiral-block}
\end{equation}
因此谱关于零对称，并且
\begin{equation}
 \det(\lambda I_{2m}-H_\tau^{(2m)}(z))
 =\det(\lambda^2I_m-B(z)^*B(z)).
 \label{eq:general-even-characteristic}
\end{equation}

两种表述从不同层次描述同一现象。式 \eqref{eq:anticommutator-coefficients} 是算子陈述：
交错对角符号反转位移一耦合，半胞平移反转位移二系数。定理~\ref{thm:flux-chiral} 则把
判据变成符号图自身的陈述：只从 switching 不变的局部圈 word $Q$ 及其半胞乘积即可读出
chirality。

这一对称把谱维数减半，却不会使任意 $m\times m$ block 可显式求解，也不保证谱边小于
$8$：周期二交错 word 已满足判据，但谱边恰为 $8$。八周期 word
\eqref{eq:target-words} 还具有两项额外性质：它首次严格越过参考谱边，而且其
$4\times4$ 平方 block 存在第二层代数对称。正是这层额外结构使完整色散可以显式求出。
